\documentclass{vvsu}

\vvsuyear{2026}

%%%%%%%%%%%%%%%%%%%

\usepackage{graphicx} % для изображений
\usepackage{tabularray} % для таблиц
\usepackage{siunitx} % для обозначений (процент, градус)
\usepackage{csquotes} % для корректной работы с кавычками
\usepackage{listings} % для листингов кода

% Список путей, где будут искаться изображения и файлы
\graphicspath{{images/}}

% Автор документа
\author{А. Е. Филатов}

% Настройка стилей для листингов кода
\input{listing_styles.tex}

%%%%%%%%%%%%%%%%%%%

\begin{document}

% Шапка
\vvsuhead{\linespread{1}\selectfont{}МИНОБРНАУКИ РОССИИ\\
\vspace{10pt}Федеральное государственное бюджетное образовательное учреждение\\
высшего образования\\
\fontsize{13}{13}\selectfont{}<<ВЛАДИВОСТОКСКИЙ ГОСУДАРСТВЕННЫЙ УНИВЕРСИТЕТ>>\\
(ФГБОУ ВО <<ВВГУ>>)\\
\vspace{10pt}\fontsize{12}{12}\selectfont{}ИНСТИТУТ ИНФОРМАЦИОННЫХ ТЕХНОЛОГИЙ И АНАЛИЗА ДАННЫХ\\
КАФЕДРА ИНФОРМАЦИОННЫХ ТЕХНОЛОГИЙ И СИСТЕМ}

% Название отчета
\title{Отчет\\По разработке системы управления автозаправочной станции на PYTHON}
\subtitle{по дисциплине\\<<Информатика и программирование>>}

% Участники работы
\member{Студент\\ гр. БИН-25-3}{А. Е. Филатов}
\member{Ассистент\\ преподавателя}{М.В. Водяницкий}

% Вывод титульника
\maketitle

% Содержание
\toc

% Глава - Введение
\section{Введение}
\label{sec:execution}
Автоматизация процессов на автозаправочных станциях (АЗС) является важным аспектом современного управления топливными ресурсами и обслуживания клиентов. В данном отчете представлена разработка системы управления АЗС на языке программирования Python. Целью данной работы является создание эффективного и надежного программного обеспечения, которое позволит оптимизировать процессы продажи топлива, учета запасов и взаимодействия с клиентами.
  
Целью данной лабороторной работы является разработка системы управления АЗС, полностью соответсвующая техническому заданию с упором на надужность данной системы и удобство использования для конечного пользователя.

  % Подглава - Принципы Проектирования
\subsection{Принципы Проектирования}
  При разработке были заложены следующие принципы проектирования:
\begin{vvsu_itemize}
    \item Модульность: Система была разработана с использованием модульного подхода, что позволяет легко добавлять новые функции и изменять существующие без влияния на другие части системы.
    \item Пользовательский интерфейс: Особое внимание было уделено созданию интуитивно понятного интерфейса для операторов АЗС, что способствует быстрому обучению и снижению ошибок при эксплуатации.
    \item Надежность: В систему были встроены механизмы обработки ошибок и резервного копирования данных для обеспечения непрерывной работы АЗС.
    \item Производительность: Оптимизация кода и использование эффективных алгоритмов позволили обеспечить высокую скорость обработки транзакций и управления запасами топлива для удобства пользования.
    \item Безопасность: Внедрение мер безопасности для защиты данных клиентов и предотвращения несанкционированного доступа к системе.
\end{vvsu_itemize}
  Данные принципы проектирования обеспечивают создание системы управления АЗС, которая отвечает современным требованиям и способствует эффективной работе автозаправочных станций.

  % Глава - Выполнение работы
\section{Выполнение работы}
\label{sec:execution}

  % Подглава - Артхитектура программы
  \subsection{Артхитектура программы}
  Архитектура представленной программы модульная и объектно-ориентированная, с чётким разделением ответственности между компонентами.
  Основные компоненты:
\begin{vvsu_itemize}
  \item Это классы, описывающие сущности АЗС и их поведение;
  \item Основной контроллер "GasStation";
  \item Система сохранения состояния "gas station data.json";
  \item Безопасность и ограничения "min level", "max level" и "can supply".
\end{vvsu_itemize}


% Подглава - Импорт библиотек
\subsection{Импорт библеотек}
  В начале программы импортируются необходимые модули и библиотеки, которые обеспечивают функциональность и поддержку различных операций.
На рисунке \ref{fig:../code_quoting/importModuls} представлен код программы.

\begin{vvsu_figure}{Листинг программы для importModuls}{fig:../code_quoting/importModuls}
  \begin{minipage}{.75\textwidth}
    \lstinputlisting[language=Python,basicstyle=\fontsize{10}{10}\linespread{1}\selectfont\ttfamily]{../code_quoting/importModuls}
  \end{minipage}
\end{vvsu_figure}

\begin{vvsu_itemize}
  \item json - для сохранения/загрузки данных в JSON-формате;
  \item os - для работы с операционной системой (очистка экрана, проверка файлов);
  \item datetime - для получения текущей даты и времени;
  \item typing - для аннотации типов (улучшает читаемость кода).
\end{vvsu_itemize}
  
  Это необходимо для обеспечения функциональности программы и взаимодействия с внешними ресурсами.


% Подглава - Класс FuelType
\subsection{Класс FuelType}
  Класс, содержащий строки с названиями видов топлива: "АИ-92", "АИ-95", "АИ-98", "ДТ".
На рисунке \ref{fig:../code_quoting/FuelType} представлен код программы.

\begin{vvsu_figure}{Листинг программы для FuelType}{fig:../code_quoting/FuelType}
  \begin{minipage}{.75\textwidth}
    \lstinputlisting[language=Python,basicstyle=\fontsize{10}{10}\linespread{1}\selectfont\ttfamily]{../code_quoting/FuelType}
  \end{minipage}
\end{vvsu_figure}
  
  Централизует символические имена топлива, уменьшает риск опечаток. Плюс данной константы, это- изменил в одном месте, то изменилось повсеместно.


% Подглава - Класс Tank
\subsection{Класс Tank}
  Моделирует резервуар для хранения топлива. Хранит информацию о типе топлива, идентификаторе, максимальном и текущем объёме, минимальном допустимом уровне и состоянии (включена/отключена).
На рисунке \ref{fig:../code_quoting/Tank} представлен код программы.

\begin{vvsu_figure}{Листинг программы для Tank}{fig:../code_quoting/Tank}
  \begin{minipage}{.75\textwidth}
    \lstinputlisting[language=Python,basicstyle=\fontsize{10}{10}\linespread{1}\selectfont\ttfamily]{../code_quoting/Tank}
  \end{minipage}
\end{vvsu_figure}

\begin{vvsu_itemize}
  \item is low() — проверяет, упал ли уровень ниже порогового;
  \item disable if low() — автоматически отключает цистерну при низком уровне;
  \item can supply(), withdraw() — обеспечивают безопасную выдачу топлива;
  \item refill() — пополнение топлива с проверкой на переполнение;
  \item to dict() и from dict() — преобразования или преобразование объекта для сохранения состояния.
\end{vvsu_itemize}
  
  Очень важный блок: модель «реальной» цистерны, с проверками безопасности (нет выдачи если мало топлива, нет переполнения при пополнении).


% Подглава - Класс Nozzle 
\subsection{Класс Nozzle}
  Связывает конкретный вид топлива с конкретной цистерной для колонки (помогает представить пистолет на колонке), а is available() говорит, можно ли из этого сопла раздавать (если связанная цистерна включена).
На рисунке \ref{fig:../code_quoting/Nozzle} представлен код программы.

\begin{vvsu_figure}{Листинг программы для Nozzle}{fig:../code_quoting/Nozzle}
  \begin{minipage}{.75\textwidth}
    \lstinputlisting[language=Python,basicstyle=\fontsize{10}{10}\linespread{1}\selectfont\ttfamily]{../code_quoting/Nozzle}
  \end{minipage}
\end{vvsu_figure}

\begin{vvsu_itemize}
  \item Связывает конкретный вид топлива с конкретной цистерной для колонки (помогает представить пистолет на колонке);
  \item is available() говорит, можно ли из этого сопла раздавать (если связанная цистерна включена).
\end{vvsu_itemize}
  
  Отделяет логику представления топлива на колонке от самой цистерны. Упрощает работу с разными соплами на одной колонке.


% Подглава - Класс Pump
\subsection{Класс Pump}
  Группирует сопла в одну колонку с идентификатором. Возвращает все сопла или только доступные (если цистерна включена).
На рисунке \ref{fig:../code_quoting/Pump} представлен код программы.

\begin{vvsu_figure}{Листинг программы для Pump}{fig:../code_quoting/Pump}
  \begin{minipage}{.75\textwidth}
    \lstinputlisting[language=Python,basicstyle=\fontsize{10}{10}\linespread{1}\selectfont\ttfamily]{../code_quoting/Pump}
  \end{minipage}
\end{vvsu_figure}

\begin{vvsu_itemize}
  \item Упрощает меню выбора колонки и отображение, какие топлива доступны;
  \item Логическая модель «колонки» соответствует реальному АЗС.
\end{vvsu_itemize}
  
  Нужен для структуры станции: колонки состоят из сопел, и на уровне Pump удобно работать с интерфейсом.


% Подглава - Класс Operation
\subsection{Класс Operation}
  Логирует каждое действие на АЗС: продажу, пополнение, аварию и т.д.
На рисунке \ref{fig:../code_quoting/Operation} представлен код программы.

\begin{vvsu_figure}{Листинг программы для Operation}{fig:../code_quoting/Operation}
  \begin{minipage}{.75\textwidth}
    \lstinputlisting[language=Python,basicstyle=\fontsize{10}{10}\linespread{1}\selectfont\ttfamily]{../code_quoting/Operation}
  \end{minipage}
\end{vvsu_figure}

  Модуль логирования для понимания, что происходило в системе.


% Подглава - Класс GasStation 
\subsection{Класс GasStation}
  Это центральный класс, объединяющий все компоненты и реализующий пользовательский интерфейс.
На рисунке \ref{fig:../code_quoting/GasStation} представлен код программы.

\begin{vvsu_figure}{Листинг программы для GasStation}{fig:../code_quoting/GasStation}
  \begin{minipage}{.75\textwidth}
    \lstinputlisting[language=Python,basicstyle=\fontsize{10}{10}\linespread{1}\selectfont\ttfamily]{../code_quoting/GasStation}
  \end{minipage}
\end{vvsu_figure}

\begin{vvsu_itemize}
  \item Инициализация (init default setup)- Создаёт 5 цистерн (включая две для АИ-95). Настраивает 8 колонок с разным набором топлива. Реализует логику распределения: колонки 14 используют первую цистерну АИ-95, 5-8- вторую.
  \item Сохранение и загрузка данных. Все данные (цистерны, операции, статистика) сохраняются в gas station data.json. При загрузке восстанавливается не только состояние, но и связи между колонками и цистернами.
  \item Основные функции меню: обслужить клиента, проверка цистерн, пополнение топлива, баланс и статистика, история операций, перекачка топлива, управление цистернами, аварийный режим.
\end{vvsu_itemize}
  
  Сердце программы: дальше идут методы, которые реализуют меню и действия.


% Подглава - Класс Pump
\subsection{Класс Pump}
  Группирует сопла в одну колонку с идентификатором. Возвращает все сопла или только доступные (если цистерна включена).
На рисунке \ref{fig:../code_quoting/Pump} представлен код программы.

\begin{vvsu_figure}{Листинг программы для Pump}{fig:../code_quoting/Pump}
  \begin{minipage}{.75\textwidth}
    \lstinputlisting[language=Python,basicstyle=\fontsize{10}{10}\linespread{1}\selectfont\ttfamily]{../code_quoting/Pump}
  \end{minipage}
\end{vvsu_figure}

\begin{vvsu_itemize}
  \item Упрощает меню выбора колонки и отображение, какие топлива доступны;
  \item Логическая модель «колонки» соответствует реальному АЗС.
\end{vvsu_itemize}

  Нужен для структуры станции: колонки состоят из сопел, и на уровне Pump удобно работать с интерфейсом.


% Подглава - Метод  init default setup
\subsection{Метод init default setup}
  Даёт реалистичную схему подключения колонок и цистерн.
На рисунке \ref{fig:../code_quoting/init_default_setup} представлен код программы.

\begin{vvsu_figure}{Листинг программы для init default setup}{fig:../code_quoting/init_default_setup}
  \begin{minipage}{.75\textwidth}
    \lstinputlisting[language=Python,basicstyle=\fontsize{10}{10}\linespread{1}\selectfont\ttfamily]{../code_quoting/init_default_setup}
  \end{minipage}
\end{vvsu_figure}

\begin{vvsu_itemize}
  \item Создает 5 цистерн разных типов;
  \item Настраивает 8 заправочных колонок;
  \item Распределяет цистерны по колонкам (особенно для АИ-95).
\end{vvsu_itemize}
  
  Важный блок для моделирования физического устройства: как именно связаны колонки и цистерны.


% Подглава - Сохранение и загрузка состояния
\subsection{Сохранение и загрузка состояния}
  Cохраняет все данные АЗС, а также восстанавливает их при запуске.
  На рисунке \ref{fig:../code_quoting/save} представлен код программы.

\begin{vvsu_figure}{Листинг программы для сохранения и загрузка состояния}{fig:../code_quoting/save}
  \begin{minipage}{.75\textwidth}
    \lstinputlisting[language=Python,basicstyle=\fontsize{10}{10}\linespread{1}\selectfont\ttfamily]{../code_quoting/save}
  \end{minipage}
\end{vvsu_figure}

\begin{vvsu_itemize}
  \item save to files - сохраняет все данные АЗС (состояние цистерн, операции, статистику) в JSON-файл;
  \item load from files - загружает данные из JSON-файла при запуске.
\end{vvsu_itemize}
  
  Надёжный способ продолжать работу с данными после перезапуска программы и обработка ошибок предотвращает падение при проблемах с файлом.


% Подглава - Вспомогательные методы
\subsection{Вспомогательные методы}
  Централизованное логирование упрощает отслеживание событий. Ограничение длины истории экономит память и делает вывод короче и понятнее.
  На рисунке \ref{fig:../code_quoting/log_operation_and_get_disabled_tanks} представлен код программы.

\begin{vvsu_figure}{Листинг программы для вспомогательных методов}{fig:../code_quoting/log_operation_and_get_disabled_tanks}
  \begin{minipage}{.75\textwidth}
    \lstinputlisting[language=Python,basicstyle=\fontsize{10}{10}\linespread{1}\selectfont\ttfamily]{../code_quoting/log_operation_and_get_disabled_tanks}
  \end{minipage}
\end{vvsu_figure}

\begin{vvsu_itemize}
  \item log operation: добавляет запись в историю и держит историю в пределах 100 различных записей.
  \item get disabled tanks: возвращает список отключённых цистерн.
\end{vvsu_itemize}
  
  Полезные утилиты для ведения истории и проверки состояния.


% Подглава - Метод serve customer
\subsection{Метод serve customer}
  Обслуживание клиента / касса, сокращённая цитата, ключевые части.
  На рисунке \ref{fig:../code_quoting/serve_customer} представлен код программы.

\begin{vvsu_figure}{Листинг программы для метода serve customer}{fig:../code_quoting/serve_customer}
  \begin{minipage}{.75\textwidth}
    \lstinputlisting[language=Python,basicstyle=\fontsize{10}{10}\linespread{1}\selectfont\ttfamily]{../code_quoting/serve_customer}
  \end{minipage}
\end{vvsu_figure}

\begin{vvsu_itemize}
  \item Ведет весь диалог продажи топлива: выбор колонки, выбор топлива, ввод литров, расчет суммы, подтверждение, списание топлива из цистерны, обновление статистики и логирование действий над заправочной станцией;
  \item Обрабатывает ошибки (недостаточно топлива, отключенная цистерна).
\end{vvsu_itemize}

  Самая важная функция для повседневной работы станции; реализует пользовательский поток продажи с базовыми проверками.

  
% Подглава -check tanks
\subsection{check tanks}
  Печатает список всех цистерн с текущими уровнями и состоянием.
  На рисунке \ref{fig:../code_quoting/check_tanks} представлен код программы.

\begin{vvsu_figure}{Листинг программы для метода check tanks}{fig:../code_quoting/check_tanks}
  \begin{minipage}{.75\textwidth}
    \lstinputlisting[language=Python,basicstyle=\fontsize{10}{10}\linespread{1}\selectfont\ttfamily]{../code_quoting/check_tanks}
  \end{minipage}
\end{vvsu_figure}

  Нужный простой дисплей статуса цистерн. Быстрый контроль состояния, удобный для оператора.


% Подглава -refill tank
\subsection{refill tank}
  Позволяет оператору добавить топливо в выбранную цистерну (с проверкой на переполнение). Логирует операцию.
  На рисунке \ref{fig:../code_quoting/refill_tank} представлен код программы.

\begin{vvsu_figure}{Листинг программы для метода refill tank}{fig:../code_quoting/refill_tank}
  \begin{minipage}{.75\textwidth}
    \lstinputlisting[language=Python,basicstyle=\fontsize{10}{10}\linespread{1}\selectfont\ttfamily]{../code_quoting/refill_tank}
  \end{minipage}
\end{vvsu_figure}

  Операция пополнения, с защитой от переполнения. Ведется учёт пополнений.


% Подглава -show balance
\subsection{show balance}
  Печатает ключевые показатели: число обслуженных машин, общий доход, сколько каждого топлива продано и сколько принесло денег.
  На рисунке \ref{fig:../code_quoting/show_balance} представлен код программы.

\begin{vvsu_figure}{Листинг программы для метода show balance}{fig:../code_quoting/show_balance}
  \begin{minipage}{.75\textwidth}
    \lstinputlisting[language=Python,basicstyle=\fontsize{10}{10}\linespread{1}\selectfont\ttfamily]{../code_quoting/show_balance}
  \end{minipage}
\end{vvsu_figure}

  Быстрое подведение итогов работы смены или дня.


% Подглава - transfer fuel
\subsection{transfer fuel}
  Позволяет перемещать топливо между цистернами того же типа (например, перенести часть АИ-95 из одной цистерны в другую). Проверяет, что есть место в приёмнике и что источник имеет достаточно топлива.
  На рисунке \ref{fig:../code_quoting/transfer_fuel} представлен код программы.

\begin{vvsu_figure}{Листинг программы для метода transfer fuel}{fig:../code_quoting/transfer_fuel}
  \begin{minipage}{.75\textwidth}
    \lstinputlisting[language=Python,basicstyle=\fontsize{10}{10}\linespread{1}\selectfont\ttfamily]{../code_quoting/transfer_fuel}
  \end{minipage}
\end{vvsu_figure}
\begin{vvsu_itemize}
  \item Практичная функция для балансировки уровней и подготовки цистерн;
  \item Помогает избежать ситуаций, когда одна цистерна пустует, а другая переполнена.
\end{vvsu_itemize}

  Реалистичная операция обслуживания склада топлива с проверками безопасности.


% Подглава - manage tanks
\subsection{manage tanks}
  Дает оператору возможность вручную включать/отключать цистерны. При включении проверяется, не ниже ли уровень минимального порога.
  На рисунке \ref{fig:../code_quoting/manage_tanks} представлен код программы.

\begin{vvsu_figure}{Листинг программы для метода manage tanks}{fig:../code_quoting/manage_tanks}
  \begin{minipage}{.75\textwidth}
    \lstinputlisting[language=Python,basicstyle=\fontsize{10}{10}\linespread{1}\selectfont\ttfamily]{../code_quoting/manage_tanks}
  \end{minipage}
\end{vvsu_figure}
\begin{vvsu_itemize}
  \item Позволяет контролировать какие цистерны используются (например, при ремонте или планировании);
  \item Предотвращает включение пустых цистерн.
\end{vvsu_itemize}

  Удобный интерфейс администрирования состояния цистерн.


% Подглава -show pumps 
\subsection{show pumps}
  Печатает для каждой колонки, какие сопла есть и доступны ли они в данный момент.
  На рисунке \ref{fig:../code_quoting/show_pumps} представлен код программы.

\begin{vvsu_figure}{Листинг программы для метода show pumps}{fig:../code_quoting/show_pumps}
  \begin{minipage}{.75\textwidth}
    \lstinputlisting[language=Python,basicstyle=\fontsize{10}{10}\linespread{1}\selectfont\ttfamily]{../code_quoting/show_pumps}
  \end{minipage}
\end{vvsu_figure}

  Наглядный статус колонок для оператора.


% Подглава - show warning
\subsection{show warning}
  При запуске меню показывает, если есть отключённые цистерны и почему (низкий уровень или отключили вручную).
  На рисунке \ref{fig:../code_quoting/emergency_and_exit_emergency} представлен код программы.

\begin{vvsu_figure}{Листинг программы для метода show warning}{fig:../code_quoting/show_warning}
  \begin{minipage}{.75\textwidth}
    \lstinputlisting[language=Python,basicstyle=\fontsize{10}{10}\linespread{1}\selectfont\ttfamily]{../code_quoting/show_warning}
  \end{minipage}
\end{vvsu_figure}

  Удобное напоминание оператору, повышает внимательность.


% Подглава - run
\subsection{run}
  Загружает данные при старте, отображает меню в цикле, читает выбор пользователя и вызывает соответствующие методы. На выходе сохраняет данные и завершает программу.
  На рисунке \ref{fig:../code_quoting/run} представлен код программы.

\begin{vvsu_figure}{Листинг программы для метода run}{fig:../code_quoting/run}
  \begin{minipage}{.75\textwidth}
    \lstinputlisting[language=Python,basicstyle=\fontsize{10}{10}\linespread{1}\selectfont\ttfamily]{../code_quoting/run}
  \end{minipage}
\end{vvsu_figure}

  Типичный простой консольный интерфейс, удобный для демонстрации и для использования вручную.


% Подглава - start program
\subsection{start program}
  Запускает программу, если файл выполнен напрямую (а не импортирован как модуль).
  На рисунке \ref{fig:../code_quoting/start_program} представлен код программы.

\begin{vvsu_figure}{Листинг программы для метода start program}{fig:../code_quoting/start_program}
  \begin{minipage}{.75\textwidth}
    \lstinputlisting[language=Python,basicstyle=\fontsize{10}{10}\linespread{1}\selectfont\ttfamily]{../code_quoting/start_program}
  \end{minipage}
\end{vvsu_figure}

  Вся программа запускается из этого блока.


% Глава - Тестирование
\section{Тестирование}
\label{sec:execution}
  Тестирование системы управления автозаправочной станцией проводилось с целью проверки корректности работы всех функциональных компонентов и обеспечения надежности системы. Были выполнены следующие виды тестирования:
\begin{vvsu_itemize}
  \item Попытка заправиться от откключенной и корректно функционирующей цистерны;
  \item Тестирование производительности: Оценка скорости обработки запросов и отклика системы при различных нагрузках.
  \item Сохранение и востоновления состояния;
  \item Ручное включение цистерн после включения;
  \item Пополнение цистерн с проверкой на переполнение.
\end{vvsu_itemize}

  Все запуски прошли успеешно, система работает корректрно и стабильно.


% Глава - Заключение
\section{Заключение}
\label{sec:execution}
  Разработный полностью работостпособный прототип системы управления автозаправочной станцией на языке Python продемонстрировал свою эффективность и надежность в ходе тестирования. Система успешно реализует основные функции, необходимые для управления АЗС, включая обслуживание клиентов, учет топлива, управление цистернами и ведение истории операций.

  Все поставленные цели были достигнуты, и система соответствует требованиям технического задания. В ходе разработки были применены принципы программирования, что обеспечило модульность, удобство использования и безопасность системы.


\end{document}
